\begin{lemma}
\label{lemma-base-change-q-integral-bottom}
With $f : X \to S$ and $n$ as in Remark \ref{remark-base-change-holds}
assume for some $q \geq 1$ we have $BC(f, n, q - 1)$. Consider
commutative diagrams
$$
\vcenter{
\xymatrix{
X \ar[d]_f &
X' \ar[d]_{f'} \ar[l] &
X'' \ar[l]^{\pi'} \ar[d]_{f''} &
Y \ar[l]^{h'} \ar[d]^e \\
S &
S' \ar[l] &
S'' \ar[l]_\pi &
T \ar[l]_{g'}
}
}
\quad\text{and}\quad
\vcenter{
\xymatrix{
X' \ar[d]_{f'} & & Y \ar[ll]^{h = h' \circ \pi'} \ar[d]^e \\
S' & & T \ar[ll]_{g = g' \circ \pi}
}
}
$$
where all squares are cartesian, $g'$ quasi-compact and quasi-separated, and
$\pi$ is integral. Let $\mathcal{F}$ be an abelian sheaf
on $T_\etale$ annihilated by $n$. If the base change map
$$
(f')^{-1}R^qg_*\mathcal{F} \longrightarrow R^qh_*e^{-1}\mathcal{F}
$$
is an isomorphism, then the base change map
$(f'')^{-1}R^qg'_*\mathcal{F} \to R^qh'_*e^{-1}\mathcal{F}$
is an isomorphism.
\end{lemma}

\begin{proof}
Since $\pi$ and $\pi'$ are integral we have $R\pi_* = \pi_*$ and
$R\pi'_* = \pi'_*$, see Lemma \ref{lemma-what-integral}.
We also have $(f')^{-1}\pi_* = \pi'_*(f'')^{-1}$. Thus we see that
$\pi'_*(f'')^{-1}R^qg'_*\mathcal{F} = (f')^{-1}R^qg_*\mathcal{F}$
and
$\pi'_*R^qh'_*e^{-1}\mathcal{F} = R^qh_*e^{-1}\mathcal{F}$.
Thus the assumption means that our map becomes an
isomorphism after applying the functor $\pi'_*$.
Hence we see that it is an isomorphism by Lemma \ref{lemma-what-integral}.
\end{proof}
